\documentclass[10pt]{article}
\usepackage{template/template}
\usepackage{tikz}
\NewsSetKeys{date={2026年11月7日},vol={025},title={每日有趣新闻},author={编辑部},note={待完成短刊｜截至北京时间11月7日上午。时政与数学已检索；数学未找到可充分核实的当日预印本。以下不是通过完整发行验收的正式版。}}
\begin{document}
\begin{documents}
\newsitem{从立案到裁定：乙叉降冰片烯反倾销调查究竟意味着什么}{国际经贸／化学材料｜11月6日｜商务部公告、路透社}{
11月6日，中国商务部宣布对原产于欧盟和日本的乙叉降冰片烯（ENB）启动反倾销调查。它是一种用于合成橡胶生产的化学原料，尤其与三元乙丙橡胶（EPDM）产业链有关。对一般读者而言，化学名称可能比关税政策还要陌生，但正是这类上游原料的价格和供应稳定性，最终影响汽车密封件、工业橡胶制品等下游行业的成本。此次调查的意义，不在于已经判定外国企业违法，而在于一场涉及价格、产业损害和贸易救济的证据程序正式开始。\par
商务部公告显示，申请人为恒河材料科技股份有限公司，申请于10月9日递交。主管部门初步审查后认为，申请人的产业代表性及其提交的证据符合立案门槛。倾销调查期为2025年7月1日至2026年6月30日；产业损害调查期则追溯至2023年1月1日。两个时期不同，是因为判断出口价格是否构成倾销，与判断国内产业是否受到损害，所需的观察窗口并不完全相同。立案以后，出口商、进口商和国内生产商可通过问卷、听证等方式提供资料，最终结论仍有待调查。\cite{mofcom2026enb,reutersenb}\par
商务部发言人在同日说明中称，申请人提供的初步材料显示，2023年至2025年有关进口数量累计增加27.35\%，2026年上半年同比增加56\%，且进口价格整体下降。这些数字来自立案申请所依据的初步证据，并不等于独立裁定已确认存在倾销及因果关系。发言人还将此案放入中欧相互贸易救济调查增加的背景，批评欧盟近期涉华调查。作为直接利害关系方的政府表态，值得记录，但不能替代对被调查企业、进口商与终端用户利益的独立评估。\cite{mofcom2026reply}\par
反倾销程序至少需要区分三件事：低价、倾销和产业损害。国际贸易中低价可能来自规模、效率、汇率或竞争策略；即使存在符合规则定义的倾销幅度，也仍需研究国内产业损害及其因果关系。贸易救济可能保护部分本土生产商，却也可能提高使用进口原料的企业成本。因此，观察后续公告时，尤其需要看调查问卷、价格比较方法、损害指标以及利害关系方的回应，而不是将立案写成已经征税。此案预计在2027年11月6日前结束，法定条件下可以延期。\cite{mofcom2026enb}
}
\newsitem{加那利群岛：短期登陆潮与长期趋势为何能够相反}{迁移与公共政策｜11月6日｜路透社}{
西班牙加那利群岛在11月前六天救起3261名海上移民，地区官员警告接待、医疗和交通能力承受严重压力。西非至加那利群岛的航线距离长、天气多变，拥挤的小船容易遭遇致命事故；本轮救援中也有死亡报告。尤其在人口仅约一万余人的耶罗岛，短期集中抵达对公共服务的影响远大于大城市。\par
但一个容易被忽视的统计事实是：尽管11月开头出现急剧增加，2026年前十个月抵达人数仍比2024年同期低近61\%。这不是自相矛盾，而是比较基期与时间尺度不同。连续几日的流量峰值可以与全年累计下降同时成立。地方政府强调即时承载能力，救援机构关注生命安全，国家层面的政策制定者则需要考虑长期趋势。三种视角回答的是不同问题，不能只选其中一种替代全部事实。\cite{reuterscanary}
}
\newsitem{美国科学技术顾问委员会会议取消}{AI与政策｜11月6日｜路透社}{
据路透社报道，美国政府取消原定11月6日举行的总统科学技术顾问委员会会议。该委员会成员包括多家大型科技企业负责人，原本拟就新设人工智能任务组提供建议。能源部一度出现新的会议日期，但随后撤回，白宫没有立即解释取消原因。会议取消本身已经得到报道，但其背后的政治或市场动机尚不能当作已证实事实；科技股近期波动及国会格局变化只能作为待检验的背景解释。\cite{reuterstech}
}
\newsitem{航天观察：Swift 的告别}{航天｜11月6日｜NASA}{
NASA在11月6日的每日天文图片栏目以《告别Swift》纪念已经结束的空间望远镜任务。Swift自2004年发射后持续观测伽马射线暴等高能天体现象，其科学价值不只在发现瞬变信号，还在于能够快速把目标位置提供给其他望远镜继续观测。新照片和观测档案应从NASA原始页面核查；本短刊没有将无法取得许可与元数据的照片复制为新闻配图。\cite{nasaswift}
}
\newsitem{数学检索说明}{数学｜11月5—7日}{
已检索arXiv近期数学论文及组合、代数、几何、数论、逻辑、概率与Lean相关关键词，但未能取得足以可靠确认作者、提交版本和主要定理的当日原始材料。本期不将旧预印本或未核查的数学主张包装成最新成果。
}
\newsitem{出版与图片说明}{编辑部}{
本短刊的贸易调查时间轴、移民趋势与航天照片尚未制作或取得合规可归档图片，亦未完成完整深稿、XeLaTeX构建及PDF视觉验收。因此仅保存已核实的部分文字，不宣称正式发行。
}
\bibliographystyle{plain}\bibliography{refs}
\end{documents}
\end{document}